\documentclass[11pt]{article}
\usepackage[margin=25mm]{geometry}
\usepackage{fontspec}
\setmainfont{Latin Modern Roman}
\usepackage{amsmath,amssymb,mathtools}
\usepackage{unicode-math}
\setmathfont{Latin Modern Math}
\usepackage{xurl,hyperref}
\hypersetup{colorlinks=true,urlcolor=blue}
\setcounter{secnumdepth}{0}
\setlength{\parindent}{0pt}
\setlength{\parskip}{5pt}
\setlength{\emergencystretch}{3em}
\newcommand{\sref}[2]{\hyperlink{source:#1:#2}{[S#2]}}
\newcommand{\cataloguescope}[1]{\paragraph{Recorded target.}#1}
\begin{document}
\section{Problems Around Polynomials — Conjecture 10}\label{statement}

\subsection{Definitions and mathematical statement}

\paragraph{Shared notation.}$\mathbb N=\{1,2,\ldots\}$, and $\mathbb Z,\mathbb Q,\mathbb R,\mathbb C$ denote the integers, rationals, reals and complex numbers. Any other conventions are specified locally.


Let $n\geq1$ and $f(z)=\sum_{j=0}^n a_jz^j$ with all $a_j>0$, and put $a_{-1}=a_{n+1}=0$. For an increasing index list $0=k_1<\cdots<k_m=n$, define its parity-change count by
\[v=\bigl|\{j:1\leq j<m,\ k_j\not\equiv k_{j+1}\pmod2\}\bigr|.\]
Real roots are counted with multiplicity.\par Put $c_j=a_j^2-a_{j-1}a_{j+1}$, and let $k_1,\ldots,k_m$ be exactly the indices for which $c_j\geq0$, in increasing order. Both endpoint indices belong to this list. Write $v_{10}(f)$ for its parity-change count.
\paragraph{Statement recorded in the supplied source.}
Is the number of real zeros of $f$ at most $v_{10}(f)$ for every such $f$? \sref{statement}{2}
\subsection{Short English statement}
Does the parity-change count among indices with nonnegative ordinary coefficient differences bound the number of real roots?
\subsection{Sources}
\begin{description}
\item[\hypertarget{source:statement:1}{S1}] ULAM.AI and the UnsolvedMath Contributors, UnsolvedMath: A Curated Collection of Open Mathematics Problems, catalogue entry. CC BY 4.0. \url{https://huggingface.co/datasets/ulamai/UnsolvedMath}.
\item[\hypertarget{source:statement:2}{S2}] Boris Shapiro, Problems Around Polynomials -- The Good, The Bad and The Ugly, Arnold Mathematical Journal 1 (2015), pp.91--104, Conjecture 10 and its complete preceding definitions. Original question and full discussion. \url{https://arxiv.org/abs/1503.05295}.

\end{description}
\label{end:statement}
\end{document}
